% =====================================================================
%  THÈME 4 — Réactif limitant et avancement
%  Niveau MOYEN — Exercices
% =====================================================================
\documentclass[11pt,a4paper]{article}
\usepackage[utf8]{inputenc}
\usepackage[T1]{fontenc}
\usepackage{lmodern}
\usepackage[french]{babel}
\usepackage{amsmath,amssymb,mathtools}
\usepackage{icomma}
\usepackage[version=4]{mhchem}
\usepackage{siunitx}
\sisetup{output-decimal-marker={,},per-mode=symbol}
\DeclareSIUnit\Molar{\mole\per\litre}
\usepackage{xcolor}
\usepackage{tikz}
\usepackage[most]{tcolorbox}
\usepackage{enumitem}
\usepackage{array}
\usepackage{fancyhdr}
\usepackage[a4paper,margin=2.1cm,headheight=26pt,headsep=8pt]{geometry}

\definecolor{primary}{RGB}{142,93,168}
\definecolor{primarydark}{RGB}{82,45,108}
\definecolor{moycol}{RGB}{198,120,9}
\newcommand{\nq}[1]{n_{\text{#1}}}
\newcommand{\Vm}{V_{\mathrm m}}

\pagestyle{fancy}\fancyhf{}
\fancyhead[L]{\small\textcolor{primarydark}{\textbf{Fiche 5} — Thème 4 : Limitant}}
\fancyhead[R]{\small\textcolor{moycol}{\textsc{Moyen}}}
\fancyfoot[C]{\small\thepage}
\renewcommand{\headrulewidth}{0.8pt}
\renewcommand{\headrule}{\hbox to\headwidth{\color{primary}\leaders\hrule height \headrulewidth\hfill}}

\newtcolorbox[auto counter]{exo}[1][]{enhanced,breakable,boxrule=1pt,arc=3pt,
  left=3mm,right=3mm,top=2.4mm,bottom=2.4mm,colback=moycol!6,colframe=moycol,
  fonttitle=\bfseries,coltitle=white,
  title={Exercice~\thetcbcounter\ \ifx\relax#1\relax\relax\else\textmd{--- \itshape #1}\fi}}

\setlength{\parindent}{0pt}
\setlength{\parskip}{3pt}

\begin{document}

\begin{center}
{\LARGE\bfseries\color{primarydark}Thème 4 — Réactif limitant}\\[1pt]
{\large\color{moycol}\textbf{Niveau Moyen}}\\[2pt]
{\itshape Objectif : limitant + tableau d'avancement + quantités finales (produits et excès). Rendement \SI{100}{\percent}.}\\
{\color{primary}\rule{0.64\linewidth}{1pt}}
\end{center}

\textbf{Données :} $\Vm=\SI{22,4}{\litre\per\mole}$ (TPN).

\begin{exo}[bilan complet par le tableau d'avancement]
Pour $\ce{2 A + 3 B -> C + 6 D}$, on part de $\nq{A,0}=\SI{0,100}{\mole}$ et $\nq{B,0}=\SI{0,200}{\mole}$.
\textbf{a)} Écris le tableau d'avancement (état initial, état $\xi$).\quad
\textbf{b)} Détermine $\xi_{\max}$ et le réactif limitant.\quad
\textbf{c)} Donne les quantités finales de \ce{C}, \ce{D} et du réactif en excès.
\end{exo}

\begin{exo}[trois réactifs]
Soit $\ce{2 K2Cr2O7 + 3 C + 8 H2SO4 -> 3 CO2 + 2 Cr2(SO4)3 + 2 K2SO4 + 8 H2O}$.
On mélange $n(\ce{K2Cr2O7})=\SI{0,10}{\mole}$, $n(\ce{C})=\SI{0,20}{\mole}$ et
$n(\ce{H2SO4})=\SI{0,30}{\mole}$.
Calcule les trois rapports $n/\nu$ et identifie le réactif limitant.
\end{exo}

\begin{exo}[zinc et acide : gaz dégagé et excès]
On fait réagir \SI{6,5}{\gram} de zinc ($M=\SI{65}{\gram\per\mole}$) avec \SI{100}{\milli\litre}
d'acide chlorhydrique à \SI{1,5}{\Molar} : $\ce{Zn + 2 HCl -> ZnCl2 + H2 ^}$.
\textbf{a)} Convertis en moles et identifie le réactif limitant.\quad
\textbf{b)} Quel volume de \ce{H2}, aux TPN, se dégage-t-il ?\quad
\textbf{c)} Quelle quantité du réactif en excès reste-t-il ?
\end{exo}

\begin{exo}[combustion incomplète en réactifs]
On brûle \SI{88}{\gram} de propane avec seulement \SI{16}{\gram} de dioxygène :
$\ce{C3H8 + 5 O2 -> 3 CO2 + 4 H2O}$.
\emph{Données :} $M(\ce{C3H8})=44$, $M(\ce{O2})=32$.
\textbf{a)} Identifie le réactif limitant.\quad
\textbf{b)} Quel volume de \ce{CO2}, aux TPN, est réellement produit ?
\end{exo}

\begin{exo}[combustion du sulfure d'hydrogène]
On brûle du sulfure d'hydrogène : $\ce{2 H2S + 3 O2 -> 2 SO2 + 2 H2O}$. On introduit
\SI{10,2}{\gram} de \ce{H2S} ($M=34$) et \SI{9,6}{\gram} de \ce{O2} ($M=32$).
\textbf{a)} Quel est le réactif limitant ?\quad
\textbf{b)} Quelle masse de \ce{SO2} ($M=64$) se forme (rendement \SI{100}{\percent}) ?
\end{exo}

\end{document}
